\documentclass[10pt]{article} 
\renewcommand{\thempfootnote}{\arabic{mpfootnote}}

\usepackage[hidelinks]{hyperref}

\usepackage[utf8]{inputenc}

\usepackage[T1]{fontenc}
\usepackage[ttdefault=true]{AnonymousPro}

\usepackage{multicol}
\setlength{\columnsep}{1cm}

\usepackage{abc}
\usepackage{graphicx}
\usepackage{musicography}

\usepackage[margin=0.5in]{geometry} 
\geometry{letterpaper}

\usepackage{fancyhdr} 
\pagestyle{fancy} % options: empty , plain , fancy
\renewcommand{\headrulewidth}{0pt} % customise the layout...
\lhead{}\chead{}\rhead{}
\lfoot{}\cfoot{\thepage}\rfoot{}

\begin{document}

\thispagestyle{empty}

\begin{abc}[name=mcleoad,options={-O= +c -F single.fmt}]
X:1
T:Miss McLeod's Reel
T:MacLeod’s Reel/Hop High Ladies (the Cake’s All Dough)/Did You Ever See the Devil Uncle Joe
L:1/8
M:C|
K:G
|: "G"GABc dBGA | B2BA BcBA | GABc dBGB | "D"A2AG AcBA | "G"GABc dBGA |
B2BA B2d2 | "C"e2ef e2ef | "D"gedB AcBA :||: "G"G2g2 edeg | B2BA BcBA | G2g2 edeg |
"D"A2AG AcBA | "G"G2g2 edeg | B2BA B2d2 | "C"e2ef e2ef | "D"gedB AcBA :|
\end{abc}
\vskip 0.25in
\noindent\makebox[\textwidth][c]{%
  \begin{minipage}{7.0in}
    \setlength{\parindent}{1.5em}

It is believed that Miss McLeod's Reel/MacLeod's Reel is originally
from Scotland and was first printed in Nathaniel Gow's \textbf{Fifth
  Collection of Strathspey Reels}
(1809).\footnote{\url{https://tunearch.org/wiki/Annotation:Miss\_McLeod's\_Reel\_(1)}
  Nathaniel Gow (1763-–1831) was a musician, composer, music
  publisher, and son of the (arguably more famous) Niel Gow. I could
  not find the \textbf{Fifth Collection of Strathspey Reels}, but many
  other Gow collections are located here:
  \url{https://imslp.org/wiki/Category:Gow,\_Nathaniel/Collections}}
By 1826, the tune was printed in America in the \textbf{New and Compleat
Preceptor for the Fife} by William Williams (Utica, NY).\footnote{\textit{Loc. cit.} I found this publication mentioned numerous times on various
websites but could not find a facsimile online.} Miss
McLeod's Reel also appears in Elias Howe's \textbf{Musician's Omnibus No.1}
(Boston, 1862, p 44) with contradance steps:\footnote{All 7 volumes can be found here: \url{https://imslp.org/wiki/Musician's\_Omnibus\_(Howe,_Elias)}}

{\parindent0pt
  \raggedright
  ~\\
First couple down the centre \\
turn half round back (lady on gent's side, gent on lady's side) \\
cast off \\
ladies chain \\
promenade 4 \\
forward and back \\
cross over to place \\
~\\
}

It is not clear when Did You Ever See the Devil Uncle Joe/Hop High
Ladies (the Cake's All Dough)\footnote{Or the many other titles: Billy
  Boy, Enterprise and Boxer, Enterprising Boxer, Gigue américaine, The
  Girl with the Handsome Face, Green Mountain, Knickerbocker Reel,
  etc.} started to become popular. The lyrics that I learned from my
southern parents in the 1990s and from my southern fellow fiddlers in
the 2000s are below but there are many (many) other lyrics as would be
expected from such a ubiquitous tune:

{\parindent0pt
  \raggedright
  ~\\
Verse: \textit{Did you ever see the devil, Uncle Joe, Uncle Joe} (3x) \\
\textit{Don't mind the weather when the wind don't blow} \\
  ~\\
Chorus: \textit{Hop high ladies, three in a row} (x3) or \textit{Hop high ladies for the cake's all dough} (x3) \\
\textit{Don't mind the weather when the wind don't blow} \\
  ~\\
\textit{Did you ever go to meetin',\footnote{\textit{Meetin'} is short for camp meeting, which is a Protestant Christian service dating from the Second Great Awakening in the early 19th century. These were originally held outside in campgrounds because of the large turnouts and lack of large buildings on the American Frontier. You can still attend camp meetings in the southeast (and other US regions) but now they are held in large, air-conditioned buildings.} Uncle Joe, Uncle Joe} (3x) \\
\textit{Don't mind the weather when the wind don't blow} \\
  ~\\
\textit{How'd ya like the weather, Uncle Joe, Uncle Joe} (3x) \\
\textit{Don't mind the weather when the wind don't blow} \\
  ~\\
}

---Kitty Steetle
  \end{minipage}
}


\end{document}
